\subsection{Conclusions}

This project investigated the improvement of a DeepSC-style semantic transceiver through baseline comparison, module-level ablation, cross-channel evaluation, and inference-strategy analysis. The same evaluation protocol was used across the main model variants so that the effects of different components could be compared under consistent conditions. The conclusions below are based on the quantitative results reported in Section 2 and the mechanism-level analysis in Section 3.

The key conclusions that emerge from the quantitative evidence are:
\begin{enumerate}
    \item \textbf{SNR-adaptive transceiver redesign is effective under AWGN low-SNR conditions.} The Improved model clearly outperforms the Baseline at \SI{-10}{dB} and \SI{-5}{dB}, with joint gains across BLEU-4 and Word Accuracy, which reduces the likelihood that the improvement is a metric artefact.
    \item \textbf{Semantic-loss regularization is not yet stable in this implementation.} The ablation study shows that removing the semantic term improves final BLEU-4 at every tested SNR point, indicating that the current fixed and SNR-scaled formulations are not well aligned with the reconstruction objective.
    \item \textbf{BERT-based confidence refinement does not produce robust net gains.} The FT-BERT refiner mitigates some severe failures compared with earlier unconstrained GPT-style corrections but still tends to over-correct when the base decoder is strong, and it is consistently beaten by the no-refinement configuration.
    \item \textbf{Beam search is the best practical enhancement in this project.} It yields consistent low-SNR gains without retraining, and its gains dominate those of LLM-based refinement in every head-to-head comparison.
    \item \textbf{Rayleigh robustness remains an open challenge.} Improvements that are effective in AWGN do not transfer directly to fading-dominant conditions, which points to a genuine gap between scalar-SNR conditioning and richer fading-state conditioning.
\end{enumerate}

From an engineering perspective, the recommended deployment configuration for the text-based semantic communication scenario studied in this project is therefore: \textbf{Improved transceiver + Beam5 decoding, without external LLM refinement}. This recommendation is grounded in the balance between effectiveness and implementation cost. The Improved transceiver introduces manageable architectural changes relative to the baseline and provides meaningful low-SNR gains. Beam search adds a modest decoding-time cost but no retraining overhead and consistently improves the most difficult cases. In contrast, external LLM refinement increases dependency complexity while offering unstable net benefit in the current setup and is therefore not recommended in its current form.

\subsection{Suggestions for Further Work}

Future work should focus not only on improving absolute performance, but also on resolving the specific limitations identified by the experiments. The results suggest that the most important next steps are richer fading-aware channel modelling, conflict-aware semantic objective design, communication-constrained post-refinement, and complexity-aware deployment evaluation. These directions are directly motivated by the observed failure modes rather than added as generic extensions.

\begin{enumerate}
    \item \textbf{Rayleigh-specific adaptation:} incorporate richer channel-state features beyond scalar SNR, such as fading statistics, short-time variance, or estimated CSI embeddings. A natural next step is to train the model with curriculum-style sampling that emphasizes deep-fade events.
    \item \textbf{Conflict-aware multi-objective training:} apply gradient balancing methods such as GradNorm-style strategies or uncertainty weighting to reduce CE--semantic objective interference \cite{kendall2018mtl, chen2018gradnorm}. Alternatively, investigate delayed activation of the semantic loss after the cross-entropy loss has converged to a reasonable level.
    \item \textbf{Constrained post-refinement:} replace free replacement with edit-distance constraints, lexical whitelist constraints, or N-best reranking over the beam output. A hybrid scheme that only accepts a replacement when it improves a calibrated confidence signal and does not decrease beam score could address the current over-correction failure.
    \item \textbf{Confidence calibration:} jointly learn decoder confidence quality so that low-confidence triggers align better with true token errors. Temperature scaling, focal-loss fine-tuning, or ensemble-based calibration are concrete technical candidates.
    \item \textbf{Domain and task expansion:} validate on additional corpora (for example WMT or news commentary) and on task-oriented semantic communication settings (question answering, intent classification) to test the generalization of the findings.
    \item \textbf{End-to-end uncertainty modelling:} jointly model channel and decoding uncertainty so that both decoding and refinement decisions can be driven by a single consistent uncertainty signal.
    \item \textbf{Beam width sweep and deployment profiling:} investigate beam widths between 3 and 10, quantify the corresponding latency and memory cost, and produce Pareto curves that inform deployment trade-offs.
\end{enumerate}

Communication-aware LLM refinement: future work should redesign LLM refinement as a communication-constrained reranking problem rather than a free editing problem. Instead of allowing BERT or another language model to directly replace tokens, the system could generate a small candidate set from beam search and then use the language model only as one component of a combined score. The final score could include decoder log probability, semantic similarity to the received latent representation, language-model fluency, and a penalty for unnecessary edits. This would preserve the benefit of pretrained linguistic knowledge while reducing the risk of semantic drift.
Channel-state-aware attention: the current SNR-adaptive attention mechanism uses scalar SNR as the main channel-quality indicator. This is sufficient for a first design but insufficient for fading-dominant channels. A future system could feed the decoder with estimated channel-state information, per-symbol reliability, or a learned channel embedding. The attention module could then assign lower weights to unreliable latent positions and higher weights to stable positions. This would make the decoder more selective under Rayleigh fading and may directly address the weak cross-channel generalization observed in this project.
Statistical significance and repeated-seed validation: the current report uses controlled single-run comparisons with identical data splits and evaluation settings. This is appropriate for an undergraduate engineering project, but future research should repeat each experiment with multiple random seeds. Mean values, standard deviations, and confidence intervals should be reported for BLEU-4, Word Accuracy, BERTScore, and sentence similarity. This would make small differences at high SNR easier to interpret and would prevent over-interpreting marginal gains.
Task-oriented evaluation: future work should test whether the reconstructed sentence is sufficient for a downstream task, such as intent classification, question answering, or command execution. This would move the system closer to goal-oriented semantic communication. In such a setting, exact sentence reconstruction may be less important than whether the receiver makes the correct task decision. A task-level metric could therefore reveal benefits that BLEU or Word Accuracy may miss.
Complexity-aware deployment study: the final recommended system uses the improved transceiver with Beam5 decoding. Although Beam5 requires no retraining, it increases inference cost relative to greedy decoding. Future work should measure latency, GPU memory usage, CPU inference time, and energy consumption. The final output should be a Pareto comparison between reconstruction quality and computational cost. This would make the project more relevant to edge-device deployment, where semantic communication is often expected to operate under limited hardware resources.
Robustness to domain shift: the current text data are sufficient for controlled comparison, but real semantic communication systems may operate on messages from different domains. Future experiments should test news text, conversational text, technical instructions, and short command-like sentences separately. This would reveal whether the learned semantic representation is general or whether it depends heavily on the training corpus distribution.
Error taxonomy construction: future work should classify decoding errors into categories such as synonym substitution, function-word error, named-entity error, tense error, word-order error, hallucinated phrase, and omission. Such an error taxonomy would make the analysis more interpretable than aggregate metrics alone. It would also help identify which errors should be handled by beam search, which by better channel modelling, and which by constrained language-model refinement.
In addition to these technical extensions, future work should also strengthen methodological breadth. Robustness claims can be reinforced through repeated runs with different random seeds and confidence-interval reporting for every metric. Such statistical reporting would make it easier to distinguish small true effects from run-to-run variance, especially when comparing close-performing variants at high SNR.

Another important direction is practical deployment profiling. While this report focuses on reconstruction quality, real deployed systems also care about latency, memory footprint, and throughput. A complete engineering evaluation should therefore include quality--cost curves, for example BLEU gain versus decoding time for greedy and beam settings, BLEU gain versus added parameters for different refinement schemes, and energy usage under repeated inference. This richer reporting would make the results directly actionable for system designers.

A longer-term research direction is to extend the semantic communication framework beyond text. Multimodal signals such as speech, images, and video carry significantly richer semantic content and are also more sensitive to channel-induced errors. Adapting the channel-aware mechanisms proposed in this project to multimodal transceivers would likely require more substantial architectural changes, but the underlying methodology, namely the hypothesis--failure--redesign loop with explicit negative-result reporting, can be reused directly.

\subsection{Final Remark}

This project shows that semantic communication transceiver design requires careful alignment between model architecture, channel condition, training objective, and decoding strategy. In the current experiments, the most reliable configuration is the Improved transceiver with Beam5 decoding and without external LLM refinement. At the same time, the mixed Rayleigh results, the semantic-loss ablation, and the FT-BERT refinement results show that several components still require further study. Future work should therefore focus on richer channel-state modelling, better-aligned semantic objectives, and more constrained refinement methods.
